\section{Future Work}
\label{sec:future-work}

Future work is prioritized around reliability, real-cluster validation, and deployment readiness.

\subsection{Priority 1: Reliability Hardening}

The latest benchmark indicates strongest performance in read/account workflows and weaker performance in \texttt{multi\_step}, \texttt{safety}, and \texttt{submission}. Immediate work should therefore focus on:

\begin{itemize}
    \item improving tool-routing consistency for multi-step workflows,
    \item reducing HITL mismatch cases in safety-sensitive tasks,
    \item improving state-transition success for submission/action flows.
\end{itemize}

These improvements should be measured with repeated full-run evaluations and pass\textasciicircum{}k analysis under the pinned local-Ollama model stack.

\subsection{Priority 2: Real Cluster Validation}

Current results are primarily based on deterministic mock-state scenarios. The next phase should validate behavior on real Slurm clusters with controlled permissions:

\begin{itemize}
    \item live scheduler variability and queue dynamics,
    \item concurrent user interactions,
    \item realistic operational failure modes.
\end{itemize}

This step is necessary to confirm external validity of the benchmark findings.

\subsection{Priority 3: Security and Governance Readiness}

Before production deployment, the system should integrate:

\begin{itemize}
    \item role-aware authentication and authorization,
    \item auditable action logs for safety-critical operations,
    \item policy controls for high-impact tool invocations.
\end{itemize}

These controls are essential for institutional HPC environments.

\subsection{Priority 4: Productization and UX}

On the frontend/runtime side, practical enhancements include:

\begin{itemize}
    \item clearer operator-facing explanations for pending confirmations,
    \item tighter per-test and per-run trace inspection for debugging,
    \item improved latency optimization in high-cost categories.
\end{itemize}

\subsection{Summary Roadmap}

In short, the near-term roadmap is: \textbf{stabilize benchmark weaknesses} $\rightarrow$ \textbf{validate on real clusters} $\rightarrow$ \textbf{harden security/governance} $\rightarrow$ \textbf{optimize deployment UX}. This sequence keeps research outcomes aligned with practical HPC adoption constraints.